% ==============================================================================
% DADOS DO TRABALHO (capa, folha de rosto e folha de aprovação)
% Único arquivo que o(a) aluno(a) precisa editar na abertura do projeto.
% ==============================================================================

\Universidade{UNIVERSIDADE NOVE DE JULHO -- UNINOVE}
\Programa{PROGRAMA DE PÓS-GRADUAÇÃO EM INFORMÁTICA E GESTÃO DO CONHECIMENTO -- PPGI}
% Para o PPGP, descomentar e ajustar também a linha \Programa acima:
% \Area{Gestão de Projetos}
% \SiglaPrograma{PPGP}

\Autor{NOME COMPLETO DO(A) AUTOR(A)} % também usado na legenda de revisores (revisao.sty)
\Titulo{TÍTULO DO TRABALHO: SUBTÍTULO (SE HOUVER)}

\Tipoexame{Pré-Qualificação} % Pré-Qualificação, Qualificação ou Defesa
\Titulacao{Mestre}           % Mestre, Mestra, Doutor ou Doutora

% Os nomes abaixo aparecem na capa, na folha de rosto, na folha de aprovação
% e na legenda de revisores do modo revisão (revisao.sty).
\Orientador{Prof. Dr. Nome do Orientador} % obrigatório
\Coorientador{Prof. Dr. Nome do Coorientador} % opcional: comentar se não houver
% \Disciplina{Seminários em Tecnologia da Informação Inteligente} % opcional (somente pré-qualificação)
\Ingresso{03/2026}                    % opcional
% \Linha{LP1 -- Modelagem e Otimização Computacional}   % opcional
\Linha{LP2 -- Sistemas Inteligentes}   % opcional
% \Linha{Gestão da Tecnologia da Informação e Conhecimento}   % opcional
\Ano{2027}

% ------------------------------------------------------------------------------
% Somente na Defesa (folha de aprovação; ativar \folhaaprovacao em main.tex)
% ------------------------------------------------------------------------------
% \DataAprovacao{dia mês e ano} % obrigatório na Defesa
% \MembroBanca{\orientador}{Universidade Nove de Julho -- Orientador}
% \MembroBanca{\coorientador}{Instituição -- Membro interno}
% \MembroBanca{Prof(a). Dr(a). Nome}{Instituição -- Membro externo}

% ------------------------------------------------------------------------------
% Rótulo exibido no carimbo do modo revisão (atualizar a cada entrega)
% ------------------------------------------------------------------------------
\versaodocumento{\tipoexame, v1}
